%% kapa: exam=1 round=36 year=2025 period=1 form=A subject=civil_law subject_ko=민법 seq=5 number=5
%% source: Q-Net 감정평가사 기출문제 (cst00309) artlSeq=5231525 file=raw/1008/5231525/2025년 제36회 감정평가사 1차시험 1교시 문제지 원본.pdf
%% conversion: kordoc 4.24.51 -> tools/gen1.py (auto); review=proofread-sonnet; flags=none
\begin{question}{5}
실종선고에 관한 설명으로 옳지 않은 것은? (다툼이 있으면 판례에 따름)
\begin{choices}
  \item 가족관계등록부상 사망한 것으로 기재되어 있는 자에 대해서는 그 사망기재의 추정력을 뒤집을 수 있는 자료가 없는 한 실종선고를 할 수 없다.
  \item 부재자가 사망할 때에 1순위의 상속인이 있는 경우, 2순위의 상속인은 특별한 사정이 없더라도 그 부재자에 대한 실종선고를 청구할 수 있다.
  \item 실종선고를 받은 자는 실종기간이 만료한 때에 사망한 것으로 본다.
  \item 실종선고의 취소가 있을 때에 실종선고를 직접원인으로 재산을 취득한 자가 선의인 경우, 특별한 사정이 없는 한 그 받은 이익이 현존하는 한도에서 반환할 의무가 있다.
  \item 실종선고의 취소가 있을 때에 실종선고를 직접원인으로 재산을 취득한 자가 악의인 경우, 특별한 사정이 없는 한 그 받은 이익에 이자를 붙여서 반환하고 손해가 있으면 이를 배상해야 한다.
\end{choices}
\end{question}
\answer{2}
